\section{Related Work}~\label{related}
We review prior work on representing state uncertainty in symbolic planning and on determining when and what to query under uncertainty.

\subsection{State Uncertainty in Symbolic Planning}~\label{rel:state_uncertainty}
We consider uncertainty about which facts hold in the current task state. Perception errors and uncertain action outcomes can leave several states consistent with the available observations. Symbolic planning represents each state as a set of facts and describes actions through preconditions and effects~\citep{fikes1971strips,mcdermott1998pddl}. Classical planning assumes a known current state and deterministic action outcomes. Probabilistic PDDL (PPDDL) extends this representation with probabilistic action effects and rewards~\citep{younes2004ppddl1}. The Relational Dynamic Influence Diagram Language (RDDL) describes probabilistic state transitions and observation models for partially observable domains~\citep{sanner2010rddl}. POMDPs represent uncertainty over possible current states through a belief and use anticipated observations in action selection~\citep{kaelbling1998planning}. Within this research direction,~\citet{zhang2015corpp} use commonsense reasoning to construct prior beliefs for POMDP planning, and \citet{sanner2010symbolic} use first order structure for symbolic dynamic programming. We use symbolic facts and action models to construct a weighted belief over reachable successor states during execution.


Maintaining an exact belief over all possible states is often computationally intractable in large state spaces. Particle representations approximate the belief with a finite set of sampled states~\citep{thrun2005probabilistic}. Online POMDP methods reduce computation by planning from the current belief without computing a policy over the entire belief space~\citep{ross2008online}. Partially Observable Monte Carlo Planning (POMCP) combines Monte Carlo belief updates with Monte Carlo tree search for online action selection~\citep{silver2010monte}. The work of~\citet{xiao2019objectsearch} extends POMCP to target object search in clutter under perception and manipulation uncertainty. 
% However, prior work on online POMDP planning focuses on action selection under state uncertainty without explicitly querying an external agent. We propose a method that integrates action selection with querying.
We ask an external agent to verify facts that distinguish the candidate states. As answers rule out inconsistent states, belief weight can concentrate on one state even while other states remain plausible. We therefore use the most probable remaining state to resume symbolic planning.


\subsection{Query Timing and Selection under Uncertainty}~\label{rel:selective_querying}
Existing methods determine query timing from uncertainty estimates, human trust, or ambiguity among candidate actions. Some methods use uncertainty in task return~\citep{singi2023hula} or implicit uncertainty estimates learned from a visual navigation policy~\citep{igbinedion2024learning}. \citet{mangalindan2025seek} model human trust in the robot as a latent POMDP state and use the trust model to select assistance requests for team performance. Online conformal prediction guides requests for expert feedback in interactive imitation learning~\citep{zhao2025conformalized}. KnowNo requests an action choice if its conformal prediction set contains multiple candidates~\citep{ren2023robots}, and IntroPlan combines conformal prediction with retrieved introspective reasoning examples~\citep{liang2024introspective}. Query timing in KnowNo and IntroPlan depends on ambiguity among candidate actions, without an explicit belief over symbolic task states.


POMDP-based methods model information requests as planning actions to determine query timing and selection through expected cumulative reward~\citep{armstrongcrews2007oracular,burks2021collaborative,burks2023harps}. Oracular POMDPs consider an oracle that reveals the exact state at a fixed cost~\citep{armstrongcrews2007oracular}. Human assisted semantic sensing incorporates human observations into the belief and selects queries through the planning objective~\citep{burks2021collaborative,burks2023harps}. Therefore, existing methods determine when and what to query based on query costs and expected rewards. We use belief confidence to decide when to query.



Query selection depends on the information needed for the task. \citet{kaipa2016addressing} use human input to resolve perception uncertainty in bin-picking. \citet{deits2013clarifying} select questions that reduce entropy in the mapping between natural language commands and entities in the environment. \citet{anand2025adaptive} use information gain to choose feedback formats at critical states for penalty function learning, considering feedback cost and availability. In contrast, our method selects the state fact that maximizes the expected reduction in belief entropy through a query to an external agent.

